% !PRINTED 353
\weberpage{353}
%<<CONT
Functionen $\nu^{\text{ter}}$ Ordnung in Bezug auf $y$, und $(4-\nu)^{\text{ter}}$ Ordnung in Bezug auf $x$. Es ist nämlich

\weq{4}{
\begin{aligned}
P_0(x,y) &= f(x) \\
P_1(x,y) &= \tfrac{1}{4}\sum_i y_i f'(x_i) \\
P_2(x,y) &= \tfrac{1}{12}\sum_i y_i y_k f''(x_i x_k) \\
P_3(x,y) &= \tfrac{1}{4}\sum x_i f'(y_i) \\
P_4(x,y) &= f(y).
\end{aligned}
}

Wenn $x$ auf der Curve $f$ liegt, so wird $P_0=0$, und eine Wurzel der Gleichung (2) verschwindet. Nehmen wir ausserdem noch $(y)$ auf der Tangente im Punkte $(x)$ an, so ist auch $P_1(x,y)=0$, und es verschwinden zwei Wurzeln von (2). Die beiden übrigen werden nach (3) durch die quadratische Gleichung
\weq{5}{6P_2+4tP_3+t^2P_4=0}
bestimmt. Wenn diese Gleichung zwei gleiche Wurzeln hat, so wird die Linie $(x,y)$ noch einen zweiten Berührungspunkt mit der Curve $f$ haben, d.\,h. sie wird Doppeltangente sein. Die Bedingung hierfür ist aber die, dass die Discriminante der Gleichung (5) verschwindet, oder dass
\weq{6}{R(x,y) = 2P_3^2 - 3P_4P_2}
gleich Null sei. Hierin ist $R(x,y)$ eine in Bezug auf jedes der beiden Systeme $(x)$ und $(y)$ homogene Function, und zwar für die $x$ vom $2^{\text{ten}}$, für die $y$ vom $6^{\text{ten}}$ Grade. Es ist aber noch zu bemerken, dass die Gleichung $R=0$ auch dann erfüllt ist, wenn $(x)$ ein beliebiger Punkt der Curve ist, und $(y)$ mit $(x)$ zusammenfällt, weil dann

\[
P_2(x,x)=P_3(x,x)=P_4(x,x)=f(x)=0
\]

wird. Sonst aber wird $R$ nur dann verschwinden, wenn die Linie $(x,y)$ eine Doppeltangente ist. Wir stellen also den Satz auf:

1. \gesperrt{Die Gleichung $R(x,y)=0$ ist, wenn $(x)$ ein Punkt der Curve $f$ und $(y)$ ein von $(x)$ verschiedener Punkt der Tangente in $(x)$ ist, die nothwendige und hinreichende Bedingung dafür, dass $(x)$ ein Berührungspunkt einer Doppeltangente sei.}

Es kommt nun darauf an, aus dieser Bedingung eine andere herzuleiten, die nur die $x$ allein enthält, und die für keine
%CONT>>
